\section{Limitations}

While the proposed framework advances the state of the art in non-intrusive SCR-ASC aging diagnostics, several
limitations stem from the modelling assumptions and the nature of the available data.

\begin{enumerate}
    \item \textbf{Reduced reaction set and lumped ASC dynamics:}
    The model retains only three of the many reactions occurring in the SCR-ASC system: the standard SCR reaction,
    ammonia adsorption/desorption, and ammonia oxidation. The fast and slow SCR reactions, NO oxidation, and other
    side reactions are neglected. Furthermore, the ammonia oxidation dynamics on the SCR and ASC catalysts are lumped
    into a single reaction with $100\%$ nitrogen selectivity, an assumption valid only above approximately
    $225\,^{\circ}\text{C}$. This lumping introduces errors in the estimated rate constants and limits the model's
    fidelity at lower temperatures or under conditions where the omitted reactions become significant.

    \item \textbf{Polynomial approximation of Arrhenius temperature dependence:}
    Temperature-dependent rate constants and the surface concentration of viable voids are approximated by
    second-order Chebyshev polynomials fitted over a bounded temperature range. As a result, the parameter estimates
    are tied to the specific temperature window of the data used for estimation and vary across test conditions (e.g.,
    RMC versus Hot-FTP). The model cannot be extrapolated beyond the fitted temperature range, and a single set of
    parameter estimates does not generalize across widely different thermal operating profiles.

    \item \textbf{Finite residence-times with in the sampling interval assumption:}
    The model derivation assumes an integer number of residence times within
    each sampling interval. When these conditions are not met exactly, the
    resulting discrepancy is absorbed into the model structure error. The
    average adsorbed ammonia concentration $\sigma(k)$ is itself an
    approximation of the true intra-sample surface dynamics, and errors in this
    approximation propagate into the $NO_x$ reduction prediction. This
    assumption breaks down when the exhaust flow rate is significantly slow.

    \item \textbf{Dependence on catalyst saturation data:}
    The aging detection framework relies on the saturated-catalyst model, which requires that a sufficient number of
    samples near catalyst saturation be present in the data. In real-world truck operation, fewer than $10\%$ of
    samples typically satisfy this condition. Drive segments that lack sufficient near-saturation data cannot be used
    for aging detection, limiting the duty cycles over which the detector is operational. The minimum sample
    requirement ($N_{sat}^{min}$) imposes a practical constraint on how quickly a detection decision can be made.

    \item \textbf{$NO_x$ sensor cross-sensitivity:}
    Commercial $NO_x$ sensors are cross-sensitive to tailpipe ammonia, introducing a non-negative directional bias into
    the tailpipe $NO_x$ measurement. Although the bounding relationship $\eta_{sat}(k) \geq \eta_y(k)$ is preserved
    under this bias, the cross-sensitivity reduces the separation between aged and degreened parameter estimates,
    thereby weakening the detector's discriminative power. The cross-sensitivity is treated as an unknown bounded
    disturbance rather than being explicitly corrected, since tailpipe ammonia measurements are unavailable on
    production trucks.

    \item \textbf{Asymptotic statistical assumptions:}
    The Wald-test detector relies on asymptotic properties of the maximum likelihood estimator: namely, that the MLE
    is normally distributed and attains the Cram\'{e}r-Rao lower bound for large sample sizes. For short drive
    segments or limited near-saturation data, these asymptotic conditions may not be satisfied, and the theoretical
    $\chi^2$ distribution of the test statistic may not hold. In such cases, the false-alarm and detection
    probabilities deviate from their analytical values.

    \item \textbf{Absence of ground-truth aging levels:}
    Validation on truck data is limited by the absence of ground-truth information on the actual degree of catalyst
    aging for individual trucks. Only the mileage difference between data collection periods serves as an indirect
    proxy for aging severity. Consequently, the detector's quantitative performance, particularly the relationship
    between the non-centrality parameter $\lambda$ and the true aging level, cannot be precisely calibrated from the
    available data.

    \item \textbf{Threshold selection with limited data:}
    The Wald-test threshold for truck data was selected empirically to balance detection and false-alarm rates on the
    available (limited) data set rather than through the Neyman-Pearson criterion, which requires sufficient data to
    characterize the test-statistic distributions under both hypotheses. A larger and more diverse data set would be
    needed to establish statistically rigorous thresholds for deployment.

    \item \textbf{Unsaturated model detector limitations on real-world data:}
    The unsaturated catalyst parameters-based detector, while effective on repeatable test-cell cycles, does not
    generalize reliably to truck data. The backward difference in its regression amplifies sensor noise, and the higher
    parameter count (nine versus three) demands stronger persistence of excitation. Its apparent success on test-cell
    data is partly an artifact of the uniform cross-sensitivity error profile across repeatable drive cycles.
\end{enumerate}